\section*{DEDICATORIA}

Al amor de mi vida, mi esposo: Homero Javier Velasteguí.
Gracias, mi amor, por ayudarme a cumplir este sueño y por acompañarme en cada paso de este camino. Gracias por tu apoyo incondicional, no solo por haber sido mi sustento durante el tiempo que tuve que vivir en Quito, mientras tú te quedabas en Esmeraldas trabajando para que a mí no me faltara nada, sino también por haber estado para mí de todas las formas en las que te necesité.
Lograr llegar hasta aquí no fue fácil. Hubo momentos de mucho estrés, cansancio, dudas y ganas de rendirme. Hubo días en los que sentía que ya no podía más, pero siempre encontré en ti la fuerza para continuar. A pesar de la distancia, nunca me sentí sola, porque siempre encontraste la manera de hacerme sentir acompañada, de escucharme, de tranquilizarme y de recordarme que sí podía lograrlo.
Fuiste mi luz al final del túnel en los momentos más oscuros, mi arcoíris después de la lluvia y ese lugar al que siempre podía volver cuando todo parecía demasiado difícil. Tu amor, tu paciencia y todos los sacrificios que hiciste por mí fueron una parte fundamental para que hoy pueda estar aquí.
La distancia solo nos demostró lo fuerte que es nuestro vínculo y todo lo que somos el uno para el otro: esposos, novios, amantes, mejores amigos y compañeros de vida. Aunque muchas veces tuvimos que estar lejos, nunca dejaste de estar presente en mi vida y en este sueño que también hiciste tuyo.
Por eso, este logro no es solamente mío. Es nuestro. Porque detrás de cada página, cada desvelo, cada momento difícil y cada paso que me llevó hasta aquí, también estuvo tu esfuerzo, tu amor y tu confianza en mí.
Gracias por creer en mí incluso cuando yo dejaba de hacerlo. Gracias por nunca dejarme sola y por darme todo lo que estuvo a tu alcance para que pudiera cumplir este sueño.
Te amo con todo mi corazón,
Eliana

